\documentclass[12pt]{article}
\usepackage[utf8]{inputenc}
\usepackage[T1]{fontenc}
\usepackage{hyperref}
\title{Walak Language Resources for K-12 Curriculum Development}
\author{Compiled by Claude Sonnet 5.5 (claude-sonnet-5-5) for Andr\'as Kornai}
\date{}
\begin{document}
\maketitle

\section{Introduction}
Walak (ISO 639-3: wlw; Glottocode wala1269) is a Trans-New Guinea language of the Dani group. It is spoken in the Baliem Valley area of Highland Papua, Indonesia. Alternative names are Lower Pyramid, Nalak and Wodo. Wikipedia gives about 20,000 speakers (2007), while Global Recordings Network gives about 1,500. The sources disagree, so both figures are unverified. See \url{https://en.wikipedia.org/wiki/Walak_language} and \url{https://www.globalrecordings.net/en/language/wlw}.

\section{Primary Educational Resources}
The searches found no primers, textbooks or school materials in Walak. Nothing indicates that the language is taught in schools.

\section{Linguistic and Reference Materials}
The searches found no dictionary or grammar. The only descriptive source found is the Wikipedia overview: \url{https://en.wikipedia.org/wiki/Walak_language}. A related page on the Walak people is at \url{https://en.wikipedia.org/wiki/Walak_people}.

\section{Religious and Community Resources}
Scripture Earth lists scripture material, including a New Testament, for Walak: \url{https://scriptureearth.org/00eng.php?iso=wlw}. Joshua Project has a people-group profile: \url{https://joshuaproject.net/people_groups/15834}. The searches found no community organisation.

\section{Digital and Audio Resources}
Global Recordings Network lists audio recordings in Walak: \url{https://www.globalrecordings.net/en/language/wlw}.

\section{Curriculum Development Context}
Very little published material was found. A curriculum would have to start from the New Testament text and the audio recordings, plus new fieldwork with the community. The scripture material is the most likely source of a written standard. Speaker numbers should be checked before any planning.
\end{document}